\subsection{Funciones aleatorias Discretas de varias variables}

Sea $\mathcal A = \set{A_i}_{i=1}^M$ un conjunto de $M$ variables aleatorias
discretas, donde cada
$A_i = \set{a_{i1}, a_{i2}, \ldots, a_{iN_i}}$ es el conjunto de eventos
elementales o valores posibles de la $i$-ésima variable, con probabilidad
$P_{A_i}(a_{ij})$.

Ahora definamos el espacio muestral conjunto $\Omega$ como el producto
cartesiano de los espacios individuales
\begin{equation}
  \Omega := A_1 \times A_2 \times \cdots \times A_M
  = \set{\vec a = (a_1, a_2, \ldots, a_M) : a_i \in A_i}
  \label{eq:omega_total}
\end{equation}

Sobre este espacio se define la probabilidad conjunta
$P(\vec a) = P(a_1, \ldots, a_M)$, la cual debe cumplir las condiciones de no
negatividad y normalización
\begin{equation}
  P(\vec a) \ge 0 \quad \forall \vec a \in \Omega
  \quad \wedge \quad
  \sum_{\vec a \in \Omega} P(\vec a) = 1
  \label{eq:prob-normalization}
\end{equation}

Si consideramos un subconjunto de variables sobre el subespacio
$\Omega_j \subseteq \Omega$, la probabilidad marginal para una configuración
$\vec a \in \Omega_j$ se obtiene sumando la probabilidad conjunta sobre las
variables que no pertenecen a $\Omega_j$, esto es
\begin{equation}
  P_{\Omega_j}(\vec a)
  = \sum_{\vec a' \notin \Omega_j} P(\vec a, \vec a')
  \label{eq:prob-marginal-omega}
\end{equation}

Sobre este espacio se define un conjunto de funciones aleatorias
$F = \set{F_j}_{j=1}^N$, donde cada $F_j: \Omega_j \to \mathbb{R}$ es una
función escalar que asigna un valor real a cada evento $\vec a \in \Omega_j$.

Usando la Delta de Kronecker $\delta_{f, F_j(\vec a)}$, la distribución de
probabilidad $P_{F_j}(f)$, i.e. la probabilidad de que $F_j$ tome el valor
específico $f$, se expresa como
\begin{equation}
  P_{F_j}(f)
  = \sum_{\vec a \in \Omega_j} \delta_{f, F_j(\vec a)} \,
  P_{\Omega_j}(\vec a)
  \label{eq:prob-fj-kronecker}
\end{equation}

Con esta distribución, la media $\bk{F_j}$ está dada por
\begin{equation}
  \bk{F_j}
  = \sum_f f \, P_{F_j}(f)
  = \sum_f f \left(
    \sum_{\vec a \in \Omega_j} \delta_{f, F_j(\vec a)} \,
    P_{\Omega_j}(\vec a)
  \right)
  \label{eq:mean-fj-delta}
\end{equation}

\begin{equation}
  \sum_f f \, \delta_{f, F_j(\vec a)} = F_j(\vec a)
  \label{eq:kronecker-filtering}
\end{equation}

Intercambiando las sumas y usando la ecuación~\ref{eq:kronecker-filtering} se
recupera la definición usual de la media:
\begin{equation}
  \bk{F_j}
  = \sum_{\vec a \in \Omega_j} F_j(\vec a) \, P_{\Omega_j}(\vec a)
  \label{eq:mean:fj}
\end{equation}

Supongamos que se quiere una nueva variable aleatoria $S$:
\begin{equation}
  S(\vec a) = \sum_{j=1}^N F_j(\vec a), \quad \vec a \in \Omega
  \label{eq:s-set}
\end{equation}

Junto a la probabilidad $P(a_1, \ldots, a_M)$, entonces la media de $S$ está
dada por

\begin{align*}
  \bk S &= \bk{\sum_{j = 1}^N F_j}
  = \sum_{a_1}\cdots \sum_{a_M} \sum_{j = 1}^N F_j(\vec{a}) P(a_1,
  \ldots, a_M)\\
  &= \sum_{j = 1}^N \sum_{a_1} \cdots \sum_{a_M} F_j(\vec{a}) P(a_1,
  \ldots, a_M)
  = \sum_{j = 1}^N \sum_{\vec{a} \in \Omega_j} F_j(\vec{a} )
  \sum_{\vec{a} \not \in \Omega_j} P(a_1, \ldots, a_M) \\
\end{align*}

Pero las sumas sobre $\vec{a} \not \in \Omega_j$ dan la probabilidad marginal
$P_{\Omega_j}(\vec{a})$ por lo que

\begin{align*}
  \bk S = \sum_{j=1}^N \sum_{\vec a \in \Omega_j} F_j(\vec a)
  P_{\Omega_j}(\vec{a})
\end{align*}

Y esta última es la media de $F_j$ por lo que la media de $S$ es igual a
la suma de las medias

\begin{equation}
  \bk S = \bk{\sum_{j = 1}^N F_j} = \sum_{j = 1}^N \bk{F_j}
  \label{eq:mean-s}
\end{equation}

Si todas las funciones aleatorias $F_j$ tienen la misma media, i.e.
$\bk{F_j} = \bk{F}$, entonces

\begin{equation}
  \bk S = \sum_{ j = 1}^N \bk{F_j}= N \bk{F}
  \label{eq:mean-s-eq}
\end{equation}

\subsection{Varianza y Desviación Estándar de funciones aleatorias
de varias variables}

Ahora para calcular la varianza $\sigma^2_S$ de $S$ se puede utilizar
la definición

\begin{align*}
  \sigma^2_S &= \bk{S^2} - \bk{S}^2 = \bk{SS} - \bk{S}^2 = \bk{
  \sum_{j=1}^N F_j \sum_{k=1}^N F_k} - \sum_{j=1}^N \bk{F_j}
  \sum_{k=1}^N \bk{F_k}
  \\
  &= \bk{\sum_{j=1}^N \sum_{k=1}^N F_jF_k} - \sum_{j=1}^N \bk{F_j}
  \sum_{k=1}^N \bk{F_k}
\end{align*}

utilizando la ecuación~\ref{eq:mean-s} esta última
igualdad se convierte en

\begin{align}
  \sigma^2_S &= \sum_{j=1}^N \sum_{k=1}^N \bk{F_jF_k} - \sum_{j=1}^N
  \bk{F_j} \sum_{k=1}^N \bk{F_k}
\end{align}

Suponiendo que $F_j$ y $F_k$ son independientes por pares, podemos hacer

\begin{align*}
  \sigma^2_S &= \sum_{j=1}^N \sum_{k\neq j}^N \bk{F_j}\bk{F_k}
  + \sum_{j=1}^N \bk{F_j^2} - \sum_{j=1}^N \bk{F_j}\sum_{k\neq j}^N \bk{F_k}
  - \sum_{j=1}^N \bk{F_j}^2
  = \sum_{j=1}^N \left(
    \bk{F_j^2} - \bk{F_j}^2
  \right)
\end{align*}

Por lo tanto tenemos que

\begin{equation}
  \sigma^2_S = \sum_{j=1}^N \sigma^2_j
  \label{eq:var-s}
\end{equation}

Con lo que podemos obtener la desviación estándar $\sigma_S$ como

\begin{equation}
  \sigma_S = \sqrt{\sigma^2_S}
  \label{eq:std-s}
\end{equation}

De la misma forma, si todas las funciones aleatorias $F_j$ tienen la misma
varianza, i.e. $\sigma^2_j = \sigma^2$, entonces

\begin{equation}
  \sigma^2_S = \sum_{j=1}^N \sigma^2_j = N \sigma^2
  \label{eq:var-s-eq}
\end{equation}

Y por lo tanto

\begin{equation}
  \sigma_S = \sigma\sqrt{N}
  \label{eq:std-s-eq}
\end{equation}

Y por lo tanto la desviación estándar relativa decrece con la raíz
cuadrada del número de variables.

\begin{equation}
  \frac{\sigma_S}{\bk{S}} = \frac{\sigma\sqrt{N}}{N\bk{F}}
  = \frac{\sigma}{\sqrt{N}\bk{F}}
  \label{eq:std-relative}
\end{equation}
